% from ega2-1-fr.tex
\oldpage[II]{14}
En effet, l'isomorphisme canonique
(\hyperref[II.1.5.2.1-fr]{1.5.2.1}) fournit un isomorphisme canonique de
$p_1^*(\widetilde{\mathcal{M}})$ sur le faisceau sur
$\operatorname{Spec}(q^*(\mathcal{B}))$ associé au
$q^*(\mathcal{B})$-Module $q^*(\mathcal{M})$, et il suffit ensuite d'appliquer
(\hyperref[II.1.4.7-fr]{1.4.7}).
\end{env}

\subsection{Morphismes affines}
\label{II.1.6-fr}

\begin{env}[1.6.1]
\label{II.1.6.1-fr}
Nous dirons qu'un morphisme $f:X\to Y$ de préschémas est \emph{affine} s'il
définit $X$ comme un préschéma affine au-dessus de $Y$. Les propriétés des
préschémas affines au-dessus d'un autre se traduisent comme suit dans ce
langage~:
\end{env}

\begin{proposition}[1.6.2]
\label{II.1.6.2-fr}
\begin{enumerate}
  \item[\rm(i)] Une immersion fermée est affine.
  \item[\rm(ii)] Le composé de deux morphismes affines est affine.
  \item[\rm(iii)] Si $f:X\to Y$ est un $S$-morphisme affine,
    $f_{(S')}:X_{(S')}\to Y_{(S')}$ est affine pour toute extension $S'\to S$
    de la base.
  \item[\rm(iv)] Si $f:X\to Y$ et $f':X'\to Y'$ sont deux $S$-morphismes
    affines,
    \[
      f\times_S f':X\times_S X'\to Y\times_S Y'
    \]
    est affine.
  \item[\rm(v)] Si $f:X\to Y$ et $g:Y\to Z$ sont deux morphismes tels que
    $g\circ f$ soit affine et $g$ séparé, alors $f$ est affine.
  \item[\rm(vi)] Si $f$ est affine, il en est de même de
    $f_{\mathrm{red}}$.
\end{enumerate}
\end{proposition}

En vertu de (\hyperref[I.5.5.12-fr]{I, 5.5.12}), il suffit de démontrer (i),
(ii) et (iii). Or (i) n'est autre que l'exemple
(\hyperref[II.1.2.2-fr]{1.2.2}), et (ii) n'est autre que
(\hyperref[II.1.3.5-fr]{1.3.5})~; enfin, (iii) découle de
(\hyperref[II.1.5.1-fr]{1.5.1}) puisque $X_{(S')}$ s'identifie au produit
$X\times_Y Y_{(S')}$ (\hyperref[I.3.3.11-fr]{I, 3.3.11}).

\begin{corollary}[1.6.3]
\label{II.1.6.3-fr}
Si $X$ est un schéma affine et $Y$ un schéma, tout morphisme $X\to Y$ est
affine.
\end{corollary}

\begin{proposition}[1.6.4]
\label{II.1.6.4-fr}
Soient $Y$ un préschéma localement noethérien, $f:X\to Y$ un morphisme de type
fini. Pour que $f$ soit affine, il faut et il suffit que $f_{\mathrm{red}}$ le
soit.
\end{proposition}

Vu (\hyperref[II.1.6.2-fr]{1.6.2, (vi)}), nous n'avons à démontrer que la
suffisance de la condition. Il suffit de prouver que si $Y$ est affine et
noethérien, $X$ est affine~; or $Y_{\mathrm{red}}$ est alors affine, donc il
en est de même de $X_{\mathrm{red}}$ par hypothèse. Or, $X$ est noethérien,
donc la conclusion résulte de (\hyperref[I.6.1.7-fr]{I, 6.1.7}).

\subsection{Fibré vectoriel associé à un Faisceau de modules}
\label{II.1.7-fr}

\begin{env}[1.7.1]
\label{II.1.7.1-fr}
Soient $A$ un anneau, $E$ un $A$-module. Rappelons qu'on appelle \emph{algèbre
symétrique} sur $E$ et qu'on note $\mathbf{S}(E)$ (ou $\mathbf{S}_A(E)$)
l'algèbre quotient de l'algèbre tensorielle $\mathbf{T}(E)$ par l'idéal
bilatère engendré par les éléments $x\otimes y-y\otimes x$, où $x,y$
parcourent $E$. L'algèbre $\mathbf{S}(E)$ est caractérisée par la propriété
universelle suivante~: si $\sigma$ est l'application canonique
$E\to\mathbf{S}(E)$ (obtenue par composition de $E\to\mathbf{T}(E)$ et de
l'application canonique $\mathbf{T}(E)\to\mathbf{S}(E)$), toute application
$A$-linéaire $E\to B$, où $B$ est une $A$-algèbre \emph{commutative}, se
factorise de façon unique en
$E\xrightarrow{\sigma}\mathbf{S}(E)\xrightarrow{g}B$, où $g$ est un
$A$-homomorphisme \emph{d'algèbres}. On déduit aussitôt de cette
caractérisation que pour deux $A$-modules $E$, $F$, on a
\[
  \mathbf{S}(E\oplus F)=\mathbf{S}(E)\otimes\mathbf{S}(F)
\]

